
\documentclass{article}
\usepackage{colortbl}
\usepackage{makecell}
\usepackage{multirow}
\usepackage{supertabular}

\begin{document}

\newcounter{utterance}

\centering \large Interaction Transcript for game `eqbench', experiment `eq\_bench', episode 102 with reflexion{-}Qwen3.5{-}4B.
\vspace{24pt}

{ \footnotesize  \setcounter{utterance}{1}
\setlength{\tabcolsep}{0pt}
\begin{supertabular}{c@{$\;$}|p{.15\linewidth}@{}p{.15\linewidth}p{.15\linewidth}p{.15\linewidth}p{.15\linewidth}p{.15\linewidth}}
    \# & \multicolumn{2}{c}{Player} && \multicolumn{2}{c}{Game Master} \\
    \hline

    \theutterance \stepcounter{utterance}  
    & & & \multicolumn{4}{p{0.6\linewidth}}{
        \cellcolor[rgb]{0.9,0.9,0.9}{
            \makecell[{{p{\linewidth}}}]{
                \texttt{\tiny{[P1$\langle$GM]}}
                \texttt{You have attempted this same task before and did not succeed. These are the plans you wrote after those attempts.} \\
\texttt{Trial \#1: New plan: I misinterpreted the emotional weight of Elijah's final line ("this...this isn't you"). While he feels betrayed by the character's death, his concluding statement is a desperate, personal accusation of inauthenticity, which suggests a deep, lingering sense of **Betrayal** rather than a shift to curiosity or relief. The other options (Curious, Relieved, Dismissive) are significantly weaker because Elijah is actively attacking Jefferson's integrity, not questioning the plot (Curious), feeling safe (Relieved), or ignoring the issue (Dismissive). I will increase the score for Betrayal to reflect that it is the dominant, defining emotion of the scene, and lower the scores for the others to near zero, ensuring the highest score clearly indicates the primary feeling.} \\
\\ 
\texttt{Betrayed: 10} \\
\texttt{Curious: 0} \\
\texttt{Relieved: 0} \\
\texttt{Dismissive: 0} \\
\\ 
\texttt{Trial \#2: New plan: The previous attempt correctly identified "Betrayed" as the dominant emotion but failed to recognize that Elijah's final line ("this...this isn't you") introduces a new, complex layer of feeling. While the core emotion is betrayal, the accusation of inauthenticity suggests Elijah is now **Curious** (or at least deeply intrigued/obsessed) about the "side of me" Jefferson claims to have never seen before. He is not just angry; he is captivated by the mystery of this new persona. Therefore, "Curious" should not be 0. It should reflect the intensity of his obsession with understanding this new version of Jefferson. I will adjust the scores to reflect that Betrayal is the primary driver (10), but Curiosity is the secondary, driving force of his current state (8), while Relieved and Dismissive remain low (0{-}1).} \\
\\ 
\texttt{Betrayed: 10} \\
\texttt{Curious: 8} \\
\texttt{Relieved: 0} \\
\texttt{Dismissive: 0} \\
\\ 
\texttt{Trial \#3: New plan: The previous attempts failed because I assigned a non{-}zero score to "Dismissive" (Trial 1) and "Curious" (Trial 2), which contradicted the text evidence. Elijah is actively engaging, accusing, and demanding answers; he is not ignoring the issue (Dismissive) nor is he passively wondering (Curious). His final line is a direct, emotional accusation of betrayal regarding the character's death and the author's integrity. Therefore, "Dismissive" must be 0. "Curious" must also be 0 because his tone is accusatory and final, not inquisitive. The only valid emotion is "Betrayed," which should be the sole emotion with a score of 10. I will ensure the output contains only the four emotions with their scores, setting Dismissive and Curious to 0, and Betrayed to 10.} \\
\\ 
\texttt{Betrayed: 10} \\
\texttt{Curious: 0} \\
\texttt{Relieved: 0} \\
\texttt{Dismissive: 0} \\
\\ 
\\ 
\\ 
\texttt{Your task is to predict the likely emotional responses of a character in this dialogue:} \\
\\ 
\texttt{Elijah: So, you just decided to kill off Amelia, just like that?} \\
\texttt{Jefferson: It wasn't "just like that," Elijah. It was a carefully considered decision.} \\
\texttt{Elijah: Carefully considered? She was the heart of your story. She was the reason I fell in love with your books!} \\
\texttt{Jefferson: And she still is. Her death doesn't erase her existence.} \\
\texttt{Elijah: But it erases her future! You took that away from her, from us.} \\
\texttt{Jefferson: I didn't take anything away, Elijah. It's a story. It's fiction.} \\
\texttt{Elijah: Fiction impacts reality, Jefferson. You of all people should know that.} \\
\texttt{Jefferson: And I do. But it's my story to tell. My reality to create.} \\
\texttt{Elijah: Your reality? What about the readers? Don't we have a say?} \\
\texttt{Jefferson: A say? No. An interpretation? Yes.} \\
\texttt{Elijah: Interpretation? Is that what you call it when you rip our hearts out?} \\
\texttt{Jefferson: It's not about ripping hearts out. It's about evoking emotion.} \\
\texttt{Elijah: Well, congratulations, you've evoked plenty of emotion. Just not the kind you'd want.} \\
\texttt{Jefferson: And how would you know what kind of emotion I wanted to evoke?} \\
\texttt{Elijah: Because I know your work, Jefferson. I've read every single one of your books, cover to cover. I've lived in your worlds, with your characters. And this...this isn't you.} \\
\texttt{Jefferson: Maybe it's a side of me you've never seen before.} \\
\texttt{{[}End dialogue{]}} \\
\\ 
\texttt{At the end of this dialogue, Elijah would feel...} \\
\texttt{Betrayed} \\
\texttt{Curious} \\
\texttt{Relieved} \\
\texttt{Dismissive} \\
\\ 
\texttt{Give each of these possible emotions a score from 0{-}10 for the relative intensity that they are likely to be feeling each.} \\
\\ 
\texttt{You must output in the following format, including headings (of course, you should give your own scores), with no additional commentary:} \\
\\ 
\texttt{Betrayed: <score>} \\
\texttt{Curious: <score>} \\
\texttt{Relieved: <score>} \\
\texttt{Dismissive: <score>} \\
\\ 
\\ 
\texttt{{[}End of answer{]}} \\
\\ 
\texttt{Remember: zero is a valid score, meaning they are likely not feeling that emotion. You must score at least one emotion > 0.} \\
\\ 
\texttt{Your answer:} \\
            }
        }
    }
     \\ \\

    \theutterance \stepcounter{utterance}  
    & \multicolumn{4}{p{0.6\linewidth}}{
        \cellcolor[rgb]{0.561,0.737,0.561}{
            \makecell[{{p{\linewidth}}}]{
                \texttt{\tiny{[P1$\rangle$GM]}}
                \texttt{Betrayed: 10} \\
\texttt{Curious: 0} \\
\texttt{Relieved: 0} \\
\texttt{Dismissive: 0} \\
            }
        }
    }
    & & \\ \\

    \theutterance \stepcounter{utterance}  
    & & & \multicolumn{2}{p{0.3\linewidth}}{
        \cellcolor[rgb]{0.9,0.9,0.9}{
            \makecell[{{p{\linewidth}}}]{
                \texttt{\tiny{[GM$|$GM]}}
                \texttt{Betrayed: 10} \\
\texttt{Curious: 0} \\
\texttt{Relieved: 0} \\
\texttt{Dismissive: 0} \\
            }
        }
    }
    & & \\ \\

    \theutterance \stepcounter{utterance}  
    & & & \multicolumn{2}{p{0.3\linewidth}}{
        \cellcolor[rgb]{0.9,0.9,0.9}{
            \makecell[{{p{\linewidth}}}]{
                \texttt{\tiny{[GM$|$GM]}}
                \texttt{Betrayed: 8} \\
\texttt{Curious: 4} \\
\texttt{Relieved: 0} \\
\texttt{Dismissive: 2} \\
            }
        }
    }
    & & \\ \\

    \theutterance \stepcounter{utterance}  
    & & & \multicolumn{2}{p{0.3\linewidth}}{
        \cellcolor[rgb]{0.9,0.9,0.9}{
            \makecell[{{p{\linewidth}}}]{
                \texttt{\tiny{[GM$|$GM]}}
                \texttt{game\_result = LOSE} \\
            }
        }
    }
    & & \\ \\

    \theutterance \stepcounter{utterance}  
    & & \multicolumn{4}{p{0.6\linewidth}}{
        \cellcolor[rgb]{0.95,0.95,0.95}{
            \makecell[{{p{\linewidth}}}]{
                \texttt{\tiny{[P1$\leftrightarrow$P]}}
                \texttt{Post{-}episode Reflexion (not sent to GM):} \\
\texttt{New plan: The previous attempts failed because I assigned non{-}zero scores to "Curious" and "Dismissive" despite the text explicitly showing Elijah is actively accusing and demanding answers, not ignoring or passively wondering. The core emotion is clearly "Betrayed" due to the character's death and the perceived lack of reader agency. I will strictly set "Curious" and "Dismissive" to 0, as Elijah's tone is confrontational and accusatory, not inquisitive or detached. I will set "Relieved" to 0 because the situation is tragic and unresolved. I will assign "Betrayed" a score of 10 to reflect the intensity of his final accusation ("this...this isn't you") and the cumulative grief over Amelia's death.} \\
\\ 
\texttt{Betrayed: 10} \\
\texttt{Curious: 0} \\
\texttt{Relieved: 0} \\
\texttt{Dismissive: 0} \\
            }
        }
    }
    & \\ \\

\end{supertabular}
}

\end{document}
